%قبل از شروع مسئله ترجیحاً چیزی ننویسید. اما اگر خواستید بنویسید، در قالب لیستی (فقرات و شمارش) بنویسید.
%برای اضافه کردن نمره، بعدی از شروع‌مسئله، [x] را اضافه کنید و به جای x  نمره مورد نظر را قرار دهید. برای زیرمسئله‌ها هم به همین صورت است.
%\شروع‌مسئله[10]
\شروع‌مسئله [15]
\متن‌سیاه{بلاک‌چین}

در این سوال شما باید سیستم ساده شده‌ای از بلاک‌چین را پیاده‌سازی کنید. ابتدا برای آشنایی با بلاک‌چین پیشنهاد می‌شود از 
\href{https://blog.goodaudience.com/blockchain-for-beginners-what-is-blockchain-519db8c6677a}{این}
لینک استفاده کنید. 
بلاک‌های شما شامل سه رشتۀ $data$، $hash$ و $nonce$ می‌باشد. $data$ رشته‌ای است که شامل اطلاعات است. $hash$ نتیجه‌ی تابع درهم‌سازی \متن‌لاتین{SHA 256} روی بلاک قبلی را ذخیره می‌کند. (در واقع نتیجه‌ی تابع درهم‌سازی \متن‌لاتین{SHA 256} روی رشتۀ $data+hash+nonce$ از بلاک قبلی را ذخیره می‌کند) برای استفاده از \متن‌لاتین{SHA 256} می‌توانید از \متن‌لاتین{hashlib} پایتون استفاده کنید. مقدار $hash$ برای بلاک اول تهی است.

 در خط اول ورودی به شما تعداد بلاک‌ها داده می‌شود. در خط بعد عدد $d$ به شما داده می‌شود که نشان می‌دهد خروجی درهم‌سازی هر بلاک باید با حداقل $d$ تا $0$ شروع شود. در هر خط از خطوط بعدی $data$ که رشته‌ای از اطلاعات است به ترتیب بلاک‌ها به عنوان ورودی داده می‌شود.
 
  ابتدا بلاک اول ساخته می‌شود و دیتای آن در آن قرار داده می‌شود. شما باید $nonce$ آن را طوری پیدا کنید که هش این بلاک با حداقل $d$ $0$ شروع شود. (به این عملیات ماینینگ گفته می‌شود) سپس به سراغ بلاک بعدی می‌روید. هش بلاک قبلی را در آن قرار می‌دهید. حال $nonce$ این بلاک را نیز طوری پیدا می‌کنید که هش آن با حداقل $d$ $0$ شروع شود. همین کار را برای تمام بلاک‌ها انجام دهید و $nonce$ ها را در هر خط از خروجی به ترتیب بلاک‌ها چاپ کنید.

%اگر مسئله چند قسمت داشت به این صورت عمل کنید:
%\شروع{زیرمسائل}
%\زیرمسئله[0] % مسئلۀ ۱آ
%\زیرمسئله % مسئلۀ ۱ب
%\زیرمسئله % مسئلۀ ۱پ
%\زیرمسئله % مسئلۀ ۱ت
%\پایان{زیرمسائل}

\solution{
	پاسخ نمونه.
}




%%%%%%%%%%%%%%%%%%%%%%%%%%%%%%%%%%%%%%%%%%%%%%%%%%%%%
% برخی دستورات مفید برای تمرین‌ها. %
%%%%%%%%%%%%%%%%%%%%%%%%%%%%%%%%%%%%%%%%%%%%%%%%%%%%%

% فرمول‌های نمونه:
%$f(x) = \Theta(x)$
%$T(x, y) \leq \log(x) + 2^y + \binom{2n}{n}$
% {\tt code\_function}


% لیست موارد بدون شماره‌گذاری:
%\شروع{فقرات}
%    \فقره مورد اول 
%        \شروع{فقرات}
%            \فقره مورد اول داخلی 
%                \شروع{فقرات}
%                    \فقره اولین مورد خیلی داخلی 
%                    \فقره دومین مورد خیلی داخلی 
%                \پایان{فقرات}
%            \فقره مورد دوم داخلی 
%        \پایان{فقرات}
%    \فقره مورد دوم 
%\پایان{فقرات}

% لیست موارد شماره‌گذاری شده:
%\شروع{شمارش}
%    \فقره یک مورد 
%    \فقره یک مورد دیگر 
%    \فقره باز هم یک مورد دیگر
%\پایان{شمارش}

% برای نوشتن یک فرمول چند خطی که یک شماره بخوره. برای حالت‌های دیگر به سایت زیر مراجعه کنید:
% https://www.overleaf.com/learn/latex/Aligning_equations_with_amsmath
%\begin{align} 
%    \begin{split}
%        (x+y)^3 &= (x+y)^2(x+y) \\
%                &= (x^2+2xy+y^2)(x+y) \\
%                &= (x^3+2x^2y+xy^2) + (x^2y+2xy^2+y^3) \\
%                &= x^3+3x^2y+3xy^2+y^3
%    \end{split}                                 
%\end{align}

% نوشتن ماتریس و آرایه دوبعدی
%\begin{align}
%    A = 
%    \begin{bmatrix}
%        A_{11} & A_{21} \\
%        A_{21} & A_{22}
%    \end{bmatrix}
%\end{align}

% اضافه کردن تصویر با شرح. برای اضافه کردن تصویر بدون شرح، فقط دستور درج‌تصویر را استفاده کنید.
%\شروع{شکل}[h!]
%\تنظیم‌ازوسط
%  \درج‌تصویر[پهنا=0.5\پهنای‌متن]{img.jpg}
%\شرح{شرح تصویر}
%\برچسب{تصویر:۱}
%\پایان{شکل}
%\پایان{زیرمسائل}

%کد زیر را ببینید.
%%\lstinputlisting{source_filename.py}
%\begin{latin}
%\begin{lstlisting}
%# Import the modules
%import sys
%import random
%
%ans = True
%
%while ans:
%question = raw_input("Ask the magic 8 ball a question: (press enter to quit) ")
%
%answers = random.randint(1,8)
%
%if question == "":
%sys.exit()
%
%elif answers == 1:
%print "It is certain"
%
%elif answers == 2:
%print "Outlook good"
%
%elif answers == 3:
%print "You may rely on it"
%
%elif answers == 4:
%print "Ask again later"
%
%elif answers == 5:
%print "Concentrate and ask again"
%
%elif answers == 6:
%print "Reply hazy, try again"
%
%elif answers == 7:
%print "My reply is no"
%
%elif answers == 8:
%print "My sources say no"
%\end{lstlisting}
%\end{latin}
